\documentclass[a4paper,oneside,12pt]{amsart}

\usepackage[spanish]{babel}
\usepackage[utf8]{inputenc}
%\usepackage[latin1]{inputenc}
\usepackage{latexsym}
\usepackage{multicol}
\usepackage[shortlabels]{enumitem}
\usepackage[heightrounded]{geometry}
\geometry{top=2cm,bottom=2cm,left=2cm,right=2cm}
\usepackage{graphicx}
\usepackage{tikz}
\usepackage{tikz-network}

\DeclareMathOperator{\cond}{cond}
\DeclareMathOperator{\Id}{Id}
\DeclareMathOperator{\re}{Re}
%\DeclareMathOperator{\im}{Im}
\newcommand{\I}{\mathbb{I}}
\newcommand{\R}{\mathbb{R}}
\newcommand{\Q}{\mathbb{Q}}
\newcommand{\C}{\mathbb{C}}
\newcommand{\Z}{\mathbb{Z}}
\newcommand{\N}{\mathbb{N}}
\DeclareMathOperator{\im}{Im}

\newcommand{\nombremateria}
{Investigaci\'on Operativa}
\newcommand{\nombrecuatrimestre}
    {Segundo cuatrimestre de 2022}
\newcommand{\numeroparcial}
    {Recuperatorio del Primer Parcial }
\newcommand{\fecha}
    {6 de Diciembre de 2022}
\newcommand{\numerotema}
    {\bf 1}

%\oddsidemargin -1cm \evensidemargin -1cm \textwidth 17.5cm
%topmargin 1.2cm \textheight 25cm
%\setlength{\textheight}{25cm}

\def \ds{\displaystyle}
\theoremstyle{definition}
\newtheorem{quest}{Ejercicio}
\thispagestyle{empty}

\begin{document}

\hfill \hfill
\begin{tabular}[c]{|c|c|c|c|}
\hline 
1 & 2 & 3 & Calificaci\'on \\ \hline \hline
\hspace{1.2 cm} & \hspace{1.2 cm} & \hspace{1.2 cm} &\hspace{1.5 cm} \\
\hspace{1.2 cm} & \hspace{1.2 cm} & \hspace{1.2 cm} &\hspace{1.5 cm} \\ \hline

\end{tabular}

\vspace{-1cm}

\noindent Nombre y apellido:

\vspace{0.2cm}

\noindent Número de libreta:

\noindent \hrulefill 

\smallskip
\renewcommand{\baselinestretch}{1.1}

\begin{center}
	\large \textbf{\nombremateria} 
	
	\numeroparcial -- \fecha
\end{center}

\vspace{.6cm}
\renewcommand{\theenumi}{\textbf{\arabic{enumi}}}

\noindent\textit{Para aprobar el parcial deberá obtenerse un puntaje mayor o igual que 4.}

%%% EMPIEZA EL EJERCICIO 1 %%%

\begin{quest} 

\verb|[2.5 pts.]|Una empresa de logística transporta granos desde la ciudad productora $p_1$ hasta la ciudad capital $p_n$. No hay una ruta que conecte directamente a ambas ciudades, sino que hay rutas que conectan ciudades intermedias $p_2, \dots, p_{n-1}$, descriptas en el conjunto $\mathcal{R}$. Es decir, $(i,j)\in\mathcal{R}$ si y sólo si existe una ruta de un solo sentido desde $p_i$ a $p_j$. La cantidad de toneladas transportadas mensualmente por una ruta puede ser una magnitud real.

Elaborar un modelo de Programación Lineal Entera Mixta que permita planificar el transporte mensual de granos desde $p_1$ hasta $p_n$, es decir, que determine cuántas toneladas deben transportarse  a lo largo de cada ruta. La empresa requiere tener en cuenta las siguientes restricciones:
\begin{enumerate}
    \item \verb|[0.25]| Para cada $(i,j)\in\mathcal{R}$ , la empresa puede transportar a lo largo de esa ruta a lo sumo $T_{ij}$ toneladas mensuales (el costo de transporte por tonelada es de $c_{ij}$ pesos).
    \item \verb|[0.25]| A $p_n$ deben llegar por lo menos $D$ toneladas de granos
    \item \verb|[1.0]| Si la empresa utiliza la ruta $(5,8)$ ($T_{58}=10$), debe abonar un impuesto fijo definido por la siguiente función continua:
    \[
    f(t) = \begin{cases}
        \frac{1}{4} t \quad \text{si $0\leq x\leq 4$} \\
        t - 3 \quad \text{si $4 < t\leq 7$} \\
        2t - 10 \quad \text{si $7 < t\leq 10$} \\
    \end{cases}
    \]
    donde t es la cantidad de toneladas transportadas mensualmente.
    \item \verb|[0.25]| Para $p_i$ ($2\leq i \leq n-1$) la cantidad de toneladas que ingresan a la ciudad debe ser igual a la cantidad saliente.
    \item \verb|[0.25]| Por $p_i$ pueden transitar a lo sumo $K_i$ toneladas por mes (el costo de tránsito es de $w_i$ pesos por tonelada).
\end{enumerate}

El objetivo es:
\begin{enumerate}
    \setcounter{enumi}{5}
    \item \verb|[0.5]| Minimizar el costo mensual.
\end{enumerate}
\end{quest}

%%% TERMINA EL EJERCICIO 1 %%%

%%% EMPIEZA EL EJERCICIO 2 %%%

\begin{quest} 

\verb|[4.5 pts.]| Se desean organizar los horarios de las exposiciones para un congreso científico internacional. Se llevará a cabo a lo largo de cinco días, simultáneamente en tres salas de la universidad. Cada día está dividido en $18$ franjas de $20$ minutos. Cada exposición dura $20$ minutos, excepto las charlas plenarias, que duran $60$ minutos. El conjunto de exposiciones viene dado por el conjunto $\mathcal{E}$ mientras que el conjunto de charlas plenarias viene dado por $\mathcal{P}$ ($\mathcal{P}\subset \mathcal{E}$). Como $|\mathcal{E}|<90$, pueden haber franjas libres.
\begin{enumerate}[A)]
   \item Teniendo en cuenta que se desea planificar en qué día, sala y franja horaria ocurrirá cada exposición, modelar linealmente las restricciones básicas del congreso:
   \begin{enumerate}[1.]
        \item \verb|[0.3]| En cada cada sala y franja horaria puede haber a lo sumo una exposición y cada exposición de $\mathcal{E}$ debe ser asignada a exactamente una franja horaria de algún día en alguna sala.
        \item \verb|[0.5]| Cada charla plenaria ocupa $3$ franjas horarias en esa sala.
        \item \verb|[0.5]| Durante una charla plenaria, no pueden haber exposiciones en las demás salas.
        \item \verb|[0.5]| Durante la tarde (franjas horarias $10$ a $18$ inclusive) del día $3$ o del día $4$ no deben haber exposiciones (disyunción excluyente)
        \item \verb|[1.0]| Salvo en el caso del ítem anterior y en el de las charlas plenarias, no puede ocurrir que las tres salas tengan una franja horaria libre simultáneamente.
        \item \verb|[0.7]| Las exposiciones del conjunto $\mathcal{E}_1\subset \mathcal{E}$ no pueden superponerse dos a dos, pero deben ser programadas para el mismo día.
   \end{enumerate}
   \item Utilizando las variables definidas en el ítem A), modelar cada una de las siguientes funciones objetivo, agregando, de ser necesario, las restricciones y variables correspondientes:
    \begin{enumerate}[1.]
        \item \verb|[0.5]| Minimizar $\ds\max_{1\leq d\leq 5} E_d$ donde $E_d$ es la cantidad total de charlas del día $d$.
        \item \verb|[0.5]| Maximizar la cantidad de exposiciones a la mañana (franjas horarias $1$ a $9$ inclusive) 
    \end{enumerate}
\end{enumerate}
\end{quest}

%%% TERMINA EL EJERCICIO 2 %%%

%%% EMPIEZA EL EJERCICIO 3 %%%


\begin{quest} \verb|[3 pts]| Labyrinthroid es un videojuego en el cual, en cada nivel, se debe llegar a la meta atravesando un laberinto y derrotando enemigos inmóviles. Tanto el mapa del nivel como la posición de los enemigos son conocidos: $\mathcal{O}$ es el conjunto de casilleros con obstáculos y $\mathcal{E}$ el conjunto de casilleros donde hay enemigos. 

Las reglas son las siguientes:
\begin{itemize}
    \item \verb|[0.5]| La protagonista puede moverse a un casillero adyacente en las direcciones $\uparrow, \rightarrow, \downarrow$ o  $\leftarrow $, siempre y cuando no haya un obstáculo (casillero gris). Cada movimiento consume $1$ unidad de tiempo.
    \item \verb|[0.25]| Moverse a un casillero donde se encuentra un enemigo y derrotarlo consume $2$ unidades de tiempo. 
    \item \verb|[0.5]| La protagonista tiene la habilidad especial de transformar un obstáculo en un casillero común, pero puede utilizar esta habilidad a lo sumo una vez por nivel. Usar la habilidad consume $3$ unidades de tiempo (desplazarse al nuevo casillero habilitado cuesta $1$ unidad de tiempo).
    \item \verb|[1.0]| Si la protagonista transita por dos casilleros consecutivos con monstruos, pierde $5$ puntos de vida. Al comenzar el nivel tiene $50$ puntos de vida. Debe llegar a la meta con al menos $5$ puntos de vida.
\end{itemize}
\begin{enumerate}[A)]
    \item Modelar linealmente las reglas de movimiento, casilleros de monstruos, habilidad especial y puntos de vida.
    \item Dado un nivel, modelar la función objetivo para cada uno de los siguientes logros:
    \begin{enumerate}[1.]
        \item \verb|[0.5]| Terminar el nivel en la menor cantidad de tiempo
        \item \verb|[0.25]| Terminar el nivel derrotando a la mayor cantidad de monstruos
    \end{enumerate}
\end{enumerate}
\textbf{Observación:} se debe resolver el problema corriendo el modelo \underline{una sola vez} en cada nivel. 


\end{quest}

%%% TERMINA EL EJERCICIO 3 %%%

%%% TERMINA EL PARCIAL %%%



\end{document}


